Prediction markets have become a large, regulated asset class. In the United States, an exchange
such as Kalshi lists hundreds of thousands of binary contracts on sports, weather, entertainment,
politics and public statements, many of them each week. Their prices are routinely read as
probabilities \citep{wolferszitzewitz2004,wolferszitzewitz2006}. A large literature asks whether
those probabilities are calibrated, documenting in particular the favourite--longshot bias
\citep{griffith1949,thalerziemba1988,snowbergwolfers2010,burgi2025}.

Much less is known about how these prices \emph{move}. Option pricing begins from a model of the
dynamics; the analogue for an event contract is largely missing. Such a model is needed to price
derivatives of prediction-market prices, to manage the risk of quoting them, to read a probability
correctly when it has just jumped, and to say what ``volatility'' even means for a security that
must end at 0 or 1.

This paper develops such a model and calibrates it on Kalshi. The starting point is simple: the
fair price of a binary contract is a bounded martingale that terminates at 0 or 1. Its total
variance over its life is therefore fixed at $\pr_0(1-\pr_0)$ (\cref{lem:variance}). What a model
must describe is \emph{when}, \emph{at which price levels} and \emph{in what shape} this fixed
variance is released.

We formalise this with a latent Gaussian score. The contract pays if the score ends above zero, so
the price is a Bachelier digital option. It solves the Black--Scholes equation with a digital payoff
and has an exact transition law (\cref{prop:m0}). The role of implied volatility is taken by an
\emph{information clock}: the rate at which the remaining uncertainty is resolved. We then add a
SABR-type stochastic volatility with a general backbone (M1) and jumps (M2). We extend the framework
to two other contract structures that Kalshi lists at scale:
\begin{itemize}
  \item \emph{ladders} of contracts on one numeric outcome, such as a daily high temperature or a
  points total. These are strips of digital options with an implied distribution, an implied
  volatility per strike and a smile.
  \item \emph{deadline} contracts that pay if an event occurs by a date. These are reduced-form
  credit claims with an implied hazard term structure and a theta.
\end{itemize}

We calibrate the three branches on several data sets:
\begin{itemize}
  \item 4.8 million hourly quote observations on 47,500 single contracts;
  \item 24,000 snapshots of 10,200 ladders in 9,200 events;
  \item 26,000 priced observations of 5,900 deadline contracts;
  \item full order-book recordings of 50 in-game sports markets.
\end{itemize}
Parameters are chosen on data before October 2026 and evaluated on October.

Our findings can be summarised as seven stylized facts.
\begin{enumerate}
  \item \textbf{Prices are fair in level and in total variance.} The variance identity holds within
  a few percent in the broad categories.
  \item \textbf{Clocks differ sharply by contract type.} In-game sports follow Stern's Brownian game
  almost exactly, with a uniform clock. Weather information is back-loaded into the final hours.
  Pre-game sports learn almost nothing. Deadline contracts carry an implied hazard that rises
  steeply into the deadline.
  \item \textbf{Variance arrives in jumps.} Half or more of price variance comes from discrete moves,
  even within games at the scale of seconds. Prices drift toward resolution more slowly than the
  Gaussian model predicts. Outside sports the one-cent grid accounts for this; in sports, jumps do.
  \item \textbf{The backbone is close to Wright--Fisher.} Price volatility scales like
  $[\pr(1-\pr)]^\beta$ with $\beta$ between 0.3 and 0.8; the logit-normal model ($\beta=1$) is
  rejected everywhere.
  \item \textbf{Volatility is stochastic but short-lived,} with half-lives of 7--21 hours outside weather.
  \item \textbf{Ladders have a smile, and it moves with the market.} The smile is sticky in moneyness
  rather than in strike. Implied widths and skews are close to calibrated, but implied tails are
  too heavy.
  \item \textbf{Tails and time are overpriced.} Implied tail mass across strikes, implied hazards
  before deadlines and implied variance before games all exceed their realized counterparts.
  Prices of deadline contracts decay toward zero at only a third to two-thirds of the model rate.
\end{enumerate}
These premia are the favourite--longshot bias in its cross-sectional and term-structure forms. The
one-cent tick and fees bound them to a few cents per contract.

The paper proceeds as follows.
\begin{itemize}
  \item \Cref{sec:lit} reviews related work.
  \item \Cref{sec:data} describes the market and the data.
  \item \Cref{sec:model} develops the model.
  \item \Cref{sec:method} describes the estimators.
  \item \Cref{sec:results} reports the results.
  \item \Cref{sec:discussion} discusses their interpretation and limitations.
  \item \Cref{sec:conclusion} concludes.
\end{itemize}
